\documentclass[11pt,letterpaper]{article}

% ============================================================
%  PACKAGES
% ============================================================
\usepackage[margin=1in]{geometry}
\usepackage{amsmath,amssymb,amsthm}
\usepackage{mathtools}
\usepackage{hyperref}
\usepackage{xcolor}
\usepackage{listings}
\usepackage{booktabs}
\usepackage{longtable}
\usepackage{array}
\usepackage{enumitem}
\usepackage{fancyhdr}
\usepackage{titlesec}
\usepackage{graphicx}
\usepackage{microtype}
\usepackage{datetime2}
\usepackage{mdframed}
\usepackage{tocloft}

% ============================================================
%  METADATA
% ============================================================
\hypersetup{
  pdftitle={PEET Drift Protocol: Prior Art Defensive Publication},
  pdfauthor={Ryan Van Gelder, Multiplicity Foundation},
  pdfsubject={Prime Entanglement Entropy Tensor Drift Measurement},
  pdfkeywords={PEET, Multiplicity Theory, PhaseMirror, drift protocol,
               prime decomposition, governance, ADR-038, ADR-039,
               ADR-040, ADR-041, Xi-Constitution},
  colorlinks=true,
  linkcolor=black,
  urlcolor=blue!70!black,
  citecolor=black
}

% ============================================================
%  THEOREM ENVIRONMENTS
% ============================================================
\theoremstyle{definition}
\newtheorem{definition}{Definition}[section]
\newtheorem{theorem}{Theorem}[section]
\newtheorem{lemma}[theorem]{Lemma}
\newtheorem{proposition}[theorem]{Proposition}
\newtheorem{corollary}[theorem]{Corollary}
\newtheorem{invariant}{Invariant}[section]
\newtheorem{axiom}{Axiom}[section]

\theoremstyle{remark}
\newtheorem{remark}{Remark}[section]

% ============================================================
%  CODE LISTINGS STYLE
% ============================================================
\definecolor{codebg}{rgb}{0.97,0.97,0.97}
\definecolor{codeframe}{rgb}{0.75,0.75,0.75}
\definecolor{kwcolor}{rgb}{0.0,0.3,0.7}
\definecolor{strcolor}{rgb}{0.2,0.5,0.2}
\definecolor{cmtcolor}{rgb}{0.5,0.5,0.5}

\lstdefinestyle{pythonstyle}{
  backgroundcolor=\color{codebg},
  frame=single,
  rulecolor=\color{codeframe},
  basicstyle=\ttfamily\small,
  keywordstyle=\color{kwcolor}\bfseries,
  stringstyle=\color{strcolor},
  commentstyle=\color{cmtcolor}\itshape,
  language=Python,
  showstringspaces=false,
  breaklines=true,
  breakatwhitespace=true,
  numbers=left,
  numberstyle=\tiny\color{cmtcolor},
  xleftmargin=1.5em,
  framexleftmargin=1.5em,
  tabsize=4,
  captionpos=b
}

\lstdefinestyle{yamlstyle}{
  backgroundcolor=\color{codebg},
  frame=single,
  rulecolor=\color{codeframe},
  basicstyle=\ttfamily\small,
  keywordstyle=\color{kwcolor},
  commentstyle=\color{cmtcolor}\itshape,
  showstringspaces=false,
  breaklines=true,
  numbers=left,
  numberstyle=\tiny\color{cmtcolor},
  xleftmargin=1.5em,
  framexleftmargin=1.5em,
  tabsize=2,
  captionpos=b
}

% ============================================================
%  HEADER / FOOTER
% ============================================================
\pagestyle{fancy}
\fancyhf{}
\rhead{\small Multiplicity Foundation --- Confidential \& Prior Art}
\lhead{\small PEET Drift Protocol v1.0}
\cfoot{\thepage}
\renewcommand{\headrulewidth}{0.4pt}

% ============================================================
%  SECTION FORMATTING
% ============================================================
\titleformat{\section}{\large\bfseries}{\thesection.}{0.5em}{}[\titlerule]
\titleformat{\subsection}{\normalsize\bfseries}{\thesubsection.}{0.5em}{}
\titleformat{\subsubsection}{\normalsize\itshape}{\thesubsubsection.}{0.5em}{}

% ============================================================
%  CUSTOM COMMANDS
% ============================================================
\newcommand{\peet}{\mathrm{PEET}}
\newcommand{\pirtm}{\textsc{PIRTM}}
\newcommand{\pmhq}{\textsc{PhaseMirror-HQ}}
\newcommand{\xicons}{\(\Xi\)-Constitution}
\newcommand{\drift}{\delta}
\newcommand{\wac}{\mathrm{WAC}}
\newcommand{\modforms}{\mathrm{ModForms}}
\newcommand{\specz}{\mathrm{Spec}\,\mathbb{Z}}
\newcommand{\norm}[1]{\left\|#1\right\|_F}

% ============================================================
%  DEFENSIVE PUBLICATION NOTICE BOX
% ============================================================
\newmdenv[
  topline=true, bottomline=true, leftline=true, rightline=true,
  linewidth=1.2pt, linecolor=red!60!black,
  backgroundcolor=red!3,
  innertopmargin=8pt, innerbottommargin=8pt,
  innerleftmargin=10pt, innerrightmargin=10pt,
  skipabove=12pt, skipbelow=12pt
]{defensivebox}

% ============================================================
%  DOCUMENT BEGIN
% ============================================================
\begin{document}

% ------------------------------------------------------------
%  TITLE PAGE
% ------------------------------------------------------------
\begin{titlepage}
\centering
\vspace*{1.5cm}

{\LARGE\bfseries PEET Drift Protocol}\\[0.4em]
{\large\bfseries Prime Entanglement Entropy Tensor\\
Drift Measurement, Governance Integration,\\
and Constitutional Hamiltonian Alignment}\\[1.2em]

\rule{0.6\textwidth}{0.8pt}\\[0.8em]

{\large \textbf{Comprehensive Technical Report \&\\
Prior Art Defensive Publication}}\\[0.6em]

{\normalsize Document Series: \texttt{MF-DR-2026-001}}\\
{\normalsize Classification: \textbf{Open / Prior Art Declaration}}\\[0.4em]

\rule{0.6\textwidth}{0.4pt}\\[1em]

\begin{tabular}{rl}
\textbf{Author:}          & Ryan Van Gelder \\
\textbf{Organization:}    & Multiplicity Foundation (501(c)(3)) \\
\textbf{Repository:}      & \texttt{MultiplicityFoundation/PhaseMirror-HQ} \\
\textbf{Governance Root:} & \texttt{73e8aadf105ee618ed9cf5...} (ledger tx\#4--5) \\
\textbf{Date:}            & April 25, 2026 \\
\textbf{Version:}         & 1.0.0 \\
\end{tabular}

\vspace{1.2em}
\rule{0.6\textwidth}{0.4pt}\\[1em]

\begin{defensivebox}
\textbf{PRIOR ART DEFENSIVE PUBLICATION NOTICE}\\[0.4em]
This document constitutes an intentional public disclosure of the
technical methods, mathematical frameworks, software architectures,
and governance protocols described herein. This disclosure is made
to establish prior art and prevent any third party from obtaining
patent protection over these concepts. All designs disclosed are
released under the \textbf{Prime Materia License} and \textbf{Apache 2.0},
as appropriate to each artifact. Publication date: \textbf{April 25, 2026}.
\end{defensivebox}

\vfill
{\small\textit{Multiplicity Foundation --- Bridgeville, Pennsylvania, USA}}
\end{titlepage}

% ------------------------------------------------------------
%  TABLE OF CONTENTS
% ------------------------------------------------------------
\tableofcontents
\newpage

% ============================================================
\section{Executive Summary}
% ============================================================

This document provides a comprehensive technical record of the
\textbf{PEET Drift Protocol} as developed within the
\pmhq{} repository under the Multiplicity Foundation.
It captures the complete analytical lineage from initial gap analysis
through mathematical formalization, governance integration design,
and the four-phase ADR development plan (ADR-038 through ADR-041).

The core problem addressed is a \textbf{structural dissonance} between
two in-house specification documents:

\begin{enumerate}[leftmargin=2em]
  \item \textbf{The Constitutional Hamiltonian} --- specifying a
        continuous Frobenius-norm drift ratio \(\delta(t)\) sampled
        at sub-millisecond cadence with autonomous kill-switch coupling.
  \item \textbf{The \xicons{}} --- specifying drift as a PEET
        entropy scalar \(\delta(n,t)\) measuring deviation from
        prime-tensor superposition, governed by exponential penalty
        and Langlands-automorphic admissibility.
\end{enumerate}

Against both specifications, the live repository state presents
a third reality: \texttt{drift: enabled: false} in
\texttt{.phase-mirror.yml}, discrete boolean MD-series rule checks
(not a continuous metric), a \texttt{maxWacAllowed: 1.05} tolerance
that is incommensurable with either entropy norm, and no named
custodian for drift response.

The resolution is a sequenced four-ADR plan that bridges
specification to implementation without breaking the Merkle-chained
governance ledger (currently at transaction \#5, Merkle root
\texttt{73e8aadf...}).

\subsection*{Key Claims for Prior Art}

The following novel combinations constitute the prior art asserted
by this publication:

\begin{enumerate}[leftmargin=2em]
  \item A prime-indexed entropy tensor \(\peet(n,t)\) used as the
        canonical unit of system drift measurement in a live
        governance-chained software repository.
  \item The coupling of PEET drift scalar output directly to a
        Merkle-chained governance ledger entry at commit time.
  \item A three-tier (Watch / Warn / Collapse) drift response
        protocol whose Collapse tier is bound to an immutable
        kill-switch file registered in the governance root.
  \item The specification of a Langlands-automorphic admissibility
        gate (\(Z_{\peet}(n) \in \modforms\)) as a formal
        precondition for state-transition lawfulness in a software
        governance system.
  \item The \textsc{WAC} (Weighted Arithmetic Contractivity)
        tolerance distinguished from and co-resident with a
        PEET entropy norm threshold in the same configuration
        artifact.
\end{enumerate}

% ============================================================
\section{Background and Context}
% ============================================================

\subsection{Multiplicity Theory}

Multiplicity Theory is a prime-indexed, recursively stable
mathematical framework wherein multiplicity --- the recurrence of
elements understood through prime number decomposition --- forms the
basis for modeling nonlinear, emergent structures across mathematics,
physics, computation, and cognition.

\begin{definition}[Multiplicity Space]
A \emph{multiplicity space} \(\mathcal{M}\) is a tuple
\((S, P, \mathcal{T}, \phi)\) where \(S\) is a state set,
\(P = \{p_1, p_2, p_3, \ldots\}\) is the ordered sequence of primes,
\(\mathcal{T}: P \times P \times \mathbb{R}_{\geq 0} \to \mathbb{R}\)
is a prime-interaction tensor, and
\(\phi: S \to \ell^2(P)\) is the prime-decomposition embedding.
\end{definition}

Rather than treating mathematical objects as static entities,
Multiplicity Theory reinterprets them as recursively generated
patterns of prime-labeled interactions, where sets and modules
become relationally governed multiplicity spaces with identity
and behavior preserved through recursive feedback loops across scales.

\subsection{PhaseMirror-HQ Repository}

\pmhq{} is the primary governance and runtime repository of the
Multiplicity Foundation. As of April 2026, it contains:

\begin{itemize}[leftmargin=2em]
  \item \pirtm{} --- Prime Interval Recursion Theory Machine runtime
  \item A Merkle-chained governance ledger (\texttt{governance/ledger.json})
        with 5 signed root commits (bootstrapped 2026-03-14)
  \item An MCP (Model Context Protocol) server with 13 immutable
        tool-boundary files
  \item ADR history through ADR-037 (Lever Branch Protocol at ADR-011,
        Tool Integrity Boundary at ADR-017)
  \item A \texttt{.phase-mirror.yml} configuration with \texttt{drift: enabled: false}
        and \texttt{maxWacAllowed: 1.05}
\end{itemize}

\subsection{The Two Specification Documents}

Two foundational documents govern the intended drift behavior:

\paragraph{The Constitutional Hamiltonian.}
Specifies a Hamiltonian operator \(H_{\text{knot}}(t)\) governing
system state evolution, with drift defined as a Frobenius-norm ratio
and an operational threshold of \(\delta < 0.10\), collapse threshold
of \(\delta > 0.17\), and a \texttt{VELOCITY=STOP} kill-switch signal.

\paragraph{The \xicons{}.}
Specifies drift as a PEET entropy deviation scalar derived from
prime-tensor superposition, with Langlands-automorphic admissibility
as the lawfulness criterion and exponential ethical penalty \(c(t) \sim e^t\)
for violation. The legislative evolution operator is:
\[
  \mathcal{L}_{t+1} = \pirtm \circ \text{CSL} \circ \text{Langlands} \circ \text{Firewall} \circ z_k
\]

% ============================================================
\section{Mathematical Framework}
% ============================================================

\subsection{Prime Entanglement Entropy Tensor}

\begin{definition}[PEET Field]
The \emph{Prime Entanglement Entropy Tensor} \(\peet(n,t)\) for
a system state \(n\) at time \(t\) is defined as:
\[
  \peet(n,t) = \sum_{p_i \mid n} \sum_{p_j \mid n}
    \mathcal{T}(p_i, p_j, t) \cdot \log\!\left(\frac{p_i p_j}{n}\right)
\]
where the outer sums range over all prime divisors of \(n\),
and \(\mathcal{T}(p_i, p_j, t)\) is the prime-interaction tensor
at time \(t\).
\end{definition}

\begin{definition}[State Superposition]
The \emph{prime-tensor superposition} of state \(n\) at time \(t\) is:
\[
  \Psi(n,t) = \sum_{p_i, p_j \in P(n)} \mathcal{T}(p_i, p_j, t) \cdot p_i p_j
\]
where \(P(n)\) denotes the multiset of prime factors of \(n\).
\end{definition}

\subsection{The PEET Drift Scalar}

\begin{definition}[PEET Drift]
\label{def:peet-drift}
The \emph{PEET drift scalar} \(\delta_{\peet}(n,t)\) is the
absolute deviation of the current state wavefunction \(\psi(n,t)\)
from its prime-tensor superposition:
\begin{equation}
  \boxed{
    \delta_{\peet}(n,t) =
    \left| \psi(n,t) - \sum_{p_i, p_j \in P(n)}
    \mathcal{T}(p_i, p_j, t) \cdot p_i p_j \right|
  }
  \label{eq:peet-drift}
\end{equation}
This quantity is dimensionless when \(\psi\) and \(\mathcal{T}\)
share the same normalization.
\end{definition}

\subsection{The Constitutional Hamiltonian Drift}

The Constitutional Hamiltonian specifies an independent drift measure
based on the Frobenius norm of operator deviation:

\begin{definition}[Hamiltonian Drift Ratio]
\label{def:ham-drift}
\begin{equation}
  \delta_H(t) = \frac{\norm{H_{\text{knot}}(t) - H_{\text{knot}}^{(0)}}}{E_{\text{free}}}
  \label{eq:ham-drift}
\end{equation}
where \(H_{\text{knot}}^{(0)}\) is the reference Hamiltonian at
system genesis, and \(E_{\text{free}}\) is the free energy normalization.
\end{definition}

\subsection{Incommensurability Theorem}

\begin{theorem}[Metric Incommensurability]
\label{thm:incommensurable}
Let \(\delta_{\peet}(n,t)\) be as in Definition~\ref{def:peet-drift}
and \(\delta_H(t)\) as in Definition~\ref{def:ham-drift}. These
metrics are dimensionally incommensurable: no scalar multiplication
\(\alpha \in \mathbb{R}_{>0}\) exists such that
\(\delta_{\peet}(n,t) = \alpha \cdot \delta_H(t)\) holds universally.
\end{theorem}

\begin{proof}[Proof Sketch]
\(\delta_{\peet}\) is an absolute deviation in the image space of
\(\psi\) (units: state-amplitude), while \(\delta_H\) is a
dimensionless energy ratio. The PEET scalar depends on prime-factor
structure of \(n\) (discrete), while \(\delta_H\) depends on
operator norm (continuous over a Hilbert space). No bijection
between \(\mathbb{Z}_{>0}\) (domain of \(\delta_{\peet}\)) and
the operator algebra (domain of \(\delta_H\)) preserves the
metric structure. \(\qed\)
\end{proof}

\begin{corollary}
The threshold \texttt{maxWacAllowed: 1.05} (a discrete rule-violation
ratio) is likewise incommensurable with both \(\delta_{\peet}\)
and \(\delta_H\). Three co-resident but mutually incomparable
drift proxies exist in the current codebase.
\end{corollary}

\subsection{Operator Contractivity}

The \pirtm{} constants module (ADR-010) defines the operator
contractivity condition:

\begin{invariant}[L0 Contractivity]
\label{inv:contractivity}
Let \(q_t = \|\Xi\| + \|\Lambda\| \cdot L_T\) be the composite
spectral radius. The system is \emph{L0-stable} if and only if:
\[
  q_t < 1 - \varepsilon
\]
where \(\varepsilon = 0.05\) (default). When \(q_t \geq 1 - \varepsilon\),
the system has exited the contraction regime and no governance
action can substitute for mathematical stabilization.
\end{invariant}

\subsection{Langlands Admissibility Gate}

The \xicons{} imposes a formal admissibility condition on all
state transitions:

\begin{definition}[Langlands Admissibility]
A state \(n \in \mathbb{N}\) is \emph{lawfully admissible} if and
only if its PEET partition function lies in the space of modular forms:
\[
  Z_{\peet}(n) \in \modforms
\]
Equivalently, the automorphic Galois representation associated to
\(n\) must admit a modular lift.
\end{definition}

\begin{definition}[PEET Partition Function]
\[
  Z_{\peet}(n) = \sum_{k \geq 1} a_k(n) \cdot q^k, \quad q = e^{2\pi i \tau}
\]
where \(a_k(n)\) are coefficients derived from the prime-indexed
tensor contractions of \(n\) at recursion depth \(k\).
\end{definition}

\subsection{Three-Tier Drift Response Model}

\begin{definition}[Drift Response Tiers]
Given constants \(0 < \varepsilon_W < \varepsilon_X < \varepsilon_C\),
define three response tiers:
\[
  \text{Tier}(n,t) = \begin{cases}
    \textsc{Watch}    & \text{if } \delta_{\peet}(n,t) \in (\varepsilon_W, \varepsilon_X] \\
    \textsc{Warn}     & \text{if } \delta_{\peet}(n,t) \in (\varepsilon_X, \varepsilon_C] \\
    \textsc{Collapse} & \text{if } \delta_{\peet}(n,t) > \varepsilon_C
  \end{cases}
\]
with proposed initial values \(\varepsilon_W = 0.05\),
\(\varepsilon_X = 0.10\), \(\varepsilon_C = \varepsilon_{\peet}\)
(to be committed as \texttt{PEET\_DRIFT\_EPSILON} in
\texttt{pirtm/constants.py}).
\end{definition}

\subsection{Exponential Ethical Penalty}

The \xicons{} imposes a time-growing penalty for sustained drift:

\begin{definition}[Ethical Drift Penalty]
For a state in the \textsc{Collapse} tier for duration \(\tau > 0\):
\[
  c(\tau) = c_0 \cdot e^{\lambda \tau}
\]
where \(c_0\) is the base penalty coefficient and \(\lambda > 0\)
is the ethical growth rate. No closed-form runtime enforcement
of this penalty currently exists in \pmhq{}.
\end{definition}

% ============================================================
\section{Repository Gap Analysis}
% ============================================================

\subsection{Current State}

The following table documents the precise gap between specification
and implementation as of the analysis date (March 25, 2026):

\begin{longtable}{p{3.2cm} p{3.8cm} p{3.8cm} p{2.5cm}}
\toprule
\textbf{Dimension} & \textbf{Specification} & \textbf{Live Repo State} & \textbf{Gap} \\
\midrule
\endfirsthead
\multicolumn{4}{l}{\small\itshape (continued)} \\
\toprule
\textbf{Dimension} & \textbf{Specification} & \textbf{Live Repo State} & \textbf{Gap} \\
\midrule
\endhead
\bottomrule
\endlastfoot

Drift enabled &
  Continuous PEET entropy scalar &
  \texttt{drift: enabled: false} &
  \textbf{Critical} \\

Metric unit &
  \(\delta_{\peet}(n,t)\) dimensionless entropy ratio &
  MD-001/002/003 boolean rules &
  \textbf{Critical} \\

Sampling cadence &
  Sub-ms (Hamiltonian) / per-commit (Xi) &
  CI/CD commit-triggered only &
  Moderate \\

Collapse threshold &
  \(\varepsilon_{\peet}\) (entropy bound) &
  \texttt{maxWacAllowed: 1.05} &
  \textbf{Critical} \\

Kill-switch binding &
  \texttt{VELOCITY=STOP} / exponential penalty &
  \texttt{rollback/kill\_switch.py} exists but unbound &
  High \\

Drift custodian &
  Named role required &
  None defined &
  High \\

Langlands gate &
  \(Z_{\peet}(n) \in \modforms\) &
  \texttt{pirtm/spectral/} scaffolding only &
  Low (future) \\

Governance ledger &
  \texttt{drift\_score} in each root commit &
  Not present in schema &
  Moderate \\
\end{longtable}

\subsection{Governance Ledger Status}

The governance ledger as of transaction \#5 (signed
\texttt{2026-03-16T01:34:20Z}, Merkle root
\texttt{73e8aadf105ee618ed9cf50af80c77a986c2b7a878704dbcca1cc68488723804})
covers 13 immutable files. The Merkle chain is intact with
\texttt{previous\_root\_tx\_id} linkage from tx\#1 through tx\#5.
No \texttt{drift\_score} field is present in any existing entry.

% ============================================================
\section{ADR Development Plan}
% ============================================================

The following four Architecture Decision Records constitute the
complete implementation roadmap. Each is a hard gate for the next.

\subsection{ADR-038: Drift Metric Definition}

\begin{mdframed}[linecolor=blue!40,linewidth=1pt,backgroundcolor=blue!2]
\textbf{Status:} Draft \quad
\textbf{Owner:} Ryan Van Gelder \quad
\textbf{Horizon:} 5 days from ratification \\
\textbf{Supersedes:} None \quad
\textbf{Required before:} ADR-039
\end{mdframed}

\paragraph{Decision.}
Adopt Equation~\eqref{eq:peet-drift} as the canonical drift scalar
for \pmhq{}. Commit \texttt{PEET\_DRIFT\_EPSILON} as a named
constant in \texttt{packages/pirtm/constants.py} adjacent to the
existing \texttt{DRIFT\_THRESHOLD}, \texttt{DEFAULT\_EPSILON},
and \texttt{OPERATOR\_CONTRACTIVITY\_BOUND} constants (ADR-010).
Implement \texttt{drift\_score(n, t) -> float} as a stub in
\texttt{pirtm/core/}.

\paragraph{Rejected alternatives.}
\begin{itemize}[leftmargin=2em,noitemsep]
  \item Re-using \texttt{DRIFT\_THRESHOLD = 0.3} --- incompatible
        units (L0-003 policy value, not PEET entropy bound).
  \item Adopting \(\delta_H(t)\) from the Constitutional Hamiltonian
        --- Theorem~\ref{thm:incommensurable} proves incommensurability.
\end{itemize}

\paragraph{Exit criterion.}
\texttt{drift\_score()} returns a \texttt{float} and passes 3
unit tests with known prime-tensor inputs. \texttt{PEET\_DRIFT\_EPSILON}
defined and referenced in at least one test.

\subsection{ADR-039: Drift Enable and Cadence Binding}

\begin{mdframed}[linecolor=blue!40,linewidth=1pt,backgroundcolor=blue!2]
\textbf{Status:} Blocked on ADR-038 \quad
\textbf{Horizon:} 10 days \\
\textbf{Required before:} ADR-040
\end{mdframed}

\paragraph{Decision.}
Activate \texttt{drift: enabled: true} in \texttt{.phase-mirror.yml}
and add \texttt{drift.metric: peet\_entropy\_delta} and
\texttt{drift.threshold: \$\{PEET\_DRIFT\_EPSILON\}}.
Wire \texttt{drift\_score()} into the \texttt{tests/} harness
for commit-triggered evaluation. Add \texttt{drift\_score: float}
to the governance ledger entry schema so each root commit records
the PEET drift at signing time.

\paragraph{Cadence decision.}
Batch-on-commit (default). Daemon heartbeat (real-time telemetry)
deferred to a future ADR; it requires the \texttt{daemon/} module
to consume the drift API, which depends on a stable interface
defined here.

\paragraph{Exit criterion.}
CI pipeline emits a \texttt{drift\_score} float on every PR.
Score appears as a named field in new governance ledger entries.
\texttt{.phase-mirror.yml} diff shows \texttt{enabled: false}
\(\to\) \texttt{enabled: true}.

\subsection{ADR-040: Drift Response and Ownership Protocol}

\begin{mdframed}[linecolor=blue!40,linewidth=1pt,backgroundcolor=blue!2]
\textbf{Status:} Blocked on ADR-039 \quad
\textbf{Horizon:} 18 days \\
\textbf{Required before:} ADR-041
\end{mdframed}

\paragraph{Decision.}
Name a \textbf{Drift Custodian} role and define the three-tier
response model from Section~3.6. Bind the Collapse tier to the
existing immutable file \texttt{rollback/kill\_switch.py}
(governance ledger hash:
\texttt{bf8aaedb61cca131f953057497d5466dce2fd5eea62298b17d02d13d41e52f4e}).
This is consistent with ADR-011's prohibition on autonomous merges.

\paragraph{Tier thresholds to commit in \texttt{contracts/system\_invariants.yaml}:}

\begin{center}
\begin{tabular}{llll}
\toprule
\textbf{Tier} & \textbf{Condition} & \textbf{Response} & \textbf{Owner} \\
\midrule
Watch    & \(\delta_{\peet} > 0.05\)              & Log only            & CI daemon \\
Warn     & \(\delta_{\peet} > 0.10\)              & Block merge, notify & Drift Custodian \\
Collapse & \(\delta_{\peet} > \varepsilon_{\peet}\) & Invoke kill-switch  & Custodian + human gate \\
\bottomrule
\end{tabular}
\end{center}

\paragraph{Exit criterion.}
ADR merged with named Custodian. Tier thresholds committed to
\texttt{contracts/system\_invariants.yaml}. Kill-switch trigger
condition documented and tested with a mock drift injection.

\subsection{ADR-041: PEET Spectral Integration (Langlands Gate)}

\begin{mdframed}[linecolor=blue!40,linewidth=1pt,backgroundcolor=blue!2]
\textbf{Status:} Blocked on ADR-040 \quad
\textbf{Horizon:} 30 days \\
\textbf{Required before:} None (terminal in sequence)
\end{mdframed}

\paragraph{Decision.}
Promote \texttt{pirtm/spectral/} from scaffolding to an active
drift signal source. Define the interface contract between
\texttt{pirtm/spectral/} and \texttt{drift\_score()} such that
the spectral module optionally supplies a PEET curvature correction
\(\kappa(n,t)\) to refine Equation~\eqref{eq:peet-drift}:
\begin{equation}
  \delta_{\peet}^{(\kappa)}(n,t) =
  \left| \psi(n,t) - \Psi(n,t) \right| \cdot \kappa(n,t)
  \label{eq:peet-drift-spectral}
\end{equation}
where \(\kappa(n,t) \in [0.98, 1.02]\) (within
\texttt{SPECTRAL\_CONSERVATION\_TOLERANCE = 0.02}).

\paragraph{Langlands binding.}
A state transition \(n \to n'\) is \emph{spectrally admissible}
only if the approximant \(\hat{Z}_{\peet}(n)\) passes the modular
form membership test. This is not required to be a full Coq proof
at ADR-041; a computable approximant is sufficient.

\paragraph{Ledger update.}
Add \texttt{pirtm/spectral/} to the immutable files list in
\texttt{governance/ledger.json} once the interface is stable,
requiring a new Merkle root (governance ledger tx\#6).

\paragraph{Exit criterion.}
\texttt{pirtm/spectral/} exports a function callable by
\texttt{drift\_score()}. One integration test exercises
the combined path. Governance ledger updated with new Merkle root.

% ============================================================
\section{Code Artifacts}
% ============================================================

\subsection{PEET Drift Scalar --- Core Implementation}

\begin{lstlisting}[
  style=pythonstyle,
  caption={%
    \texttt{pirtm/core/drift.py} --- PEET drift scalar (ADR-038 deliverable)
  },
  label={lst:drift-core}
]
"""
ADR-038: PEET Drift Scalar
Canonical implementation of the prime-tensor entropy deviation metric.

Units: dimensionless (normalized by state amplitude)
DO NOT modify thresholds without ADR approval.
"""

from __future__ import annotations
import math
from typing import Callable

from pirtm.constants import (
    PEET_DRIFT_EPSILON,
    DEFAULT_EPSILON,
    SPECTRAL_CONSERVATION_TOLERANCE,
)
from pirtm.governance.telemetry import telemetry


def prime_factors(n: int) -> list[int]:
    """Return sorted list of prime factors of n (with multiplicity)."""
    factors: list[int] = []
    d = 2
    while d * d <= n:
        while n % d == 0:
            factors.append(d)
            n //= d
        d += 1
    if n > 1:
        factors.append(n)
    return factors


def peet_superposition(
    n: int,
    tensor_fn: Callable[[int, int, float], float],
    t: float,
) -> float:
    """
    Compute the prime-tensor superposition Psi(n, t).

    Psi(n, t) = sum_{p_i, p_j in P(n)} T(p_i, p_j, t) * p_i * p_j

    Args:
        n:         State integer whose prime structure is analyzed.
        tensor_fn: T(p_i, p_j, t) -> float, prime-interaction tensor.
        t:         Time parameter.

    Returns:
        float: The superposition value.
    """
    primes = prime_factors(n)
    acc = 0.0
    for pi in primes:
        for pj in primes:
            acc += tensor_fn(pi, pj, t) * pi * pj
    return acc


def drift_score(
    n: int,
    psi: float,
    tensor_fn: Callable[[int, int, float], float],
    t: float,
    kappa: float = 1.0,
) -> float:
    """
    Compute the PEET drift scalar delta(n, t).

    delta(n, t) = |psi(n, t) - Psi(n, t)| * kappa

    where kappa is the optional spectral curvature correction
    from pirtm/spectral/ (ADR-041). Default kappa=1.0 (no correction).

    Args:
        n:         State integer.
        psi:       Current wavefunction amplitude psi(n, t).
        tensor_fn: Prime-interaction tensor T(p_i, p_j, t).
        t:         Time parameter.
        kappa:     Spectral curvature correction in [0.98, 1.02].

    Returns:
        float: Non-negative PEET drift scalar.

    Raises:
        ValueError: If n < 1 or kappa is outside tolerance band.
    """
    if n < 1:
        raise ValueError(f"State integer n must be >= 1, got {n}")

    tol = SPECTRAL_CONSERVATION_TOLERANCE
    if not (1.0 - tol <= kappa <= 1.0 + tol):
        raise ValueError(
            f"kappa={kappa} outside spectral tolerance "
            f"[{1.0 - tol}, {1.0 + tol}]"
        )

    superposition = peet_superposition(n, tensor_fn, t)
    raw_delta = abs(psi - superposition) * kappa

    telemetry.gauge("pirtm_drift_score", raw_delta,
                    labels={"n": str(n), "t": f"{t:.3f}"})

    return raw_delta


def classify_drift_tier(delta: float) -> str:
    """
    Classify drift scalar into governance tier (ADR-040).

    Returns one of: 'watch', 'warn', 'collapse', 'nominal'.
    """
    if delta > PEET_DRIFT_EPSILON:
        tier = "collapse"
    elif delta > 0.10:
        tier = "warn"
    elif delta > 0.05:
        tier = "watch"
    else:
        tier = "nominal"

    telemetry.counter(
        "pirtm_drift_tier_total",
        labels={"tier": tier}
    )
    return tier
\end{lstlisting}

\subsection{Configuration Activation (ADR-039)}

\begin{lstlisting}[
  style=yamlstyle,
  caption={\texttt{.phase-mirror.yml} --- proposed diff for ADR-039},
  label={lst:phase-mirror-yml}
]
version: '1'
rules:
  enabled:
    - MD-001
    - MD-002
    - MD-003
  severity:
    MD-001: critical
    MD-002: high
    MD-003: high
l0_invariants:
  enabled: true
  strict: false
drift:
  enabled: true                    # ADR-039: was false
  metric: peet_entropy_delta       # ADR-039: PEET scalar (eq. 1)
  threshold: 0.30                  # ADR-038: PEET_DRIFT_EPSILON
  tiers:
    watch:    0.05
    warn:     0.10
    collapse: 0.30                 # binds to rollback/kill_switch.py
  cadence: commit                  # batch-on-commit (daemon deferred)
  custodian: "@rvangelder"         # ADR-040: named Drift Custodian

levers:
  enabled: false
  maxConcurrentBranches: 5
  defaultTension: 0.5
  maxWacAllowed: 1.05
  autoMergeIfCertified: false
  mergeStrategy: squash
\end{lstlisting}

\subsection{Governance Ledger Schema Extension (ADR-039)}

\begin{lstlisting}[
  style=pythonstyle,
  caption={Proposed ledger entry schema extension --- \texttt{governance/ledger.py}},
  label={lst:ledger-schema}
]
# ADR-039: Add drift_score field to governance ledger entries.
# Each GOVERNANCE_ROOT_COMMIT must record the PEET drift scalar
# computed at time of signing.

LEDGER_ENTRY_SCHEMA_V2 = {
    "type": "GOVERNANCE_ROOT_COMMIT",
    "governance_version": str,
    "merkle_root": str,
    "drift_score": float,          # ADR-039: PEET delta at commit time
    "drift_tier": str,             # ADR-040: 'nominal'|'watch'|'warn'
    "immutable_files": list,
    "signed_by": str,
    "timestamp": str,
    "previous_root_tx_id": int,
    "governance_notes": str,
}

def new_root_commit(
    merkle_root: str,
    immutable_files: list,
    signed_by: str,
    drift_score: float,
    previous_tx_id: int,
    notes: str = "",
) -> dict:
    """Construct a v2 governance ledger entry with drift tracking."""
    from pirtm.core.drift import classify_drift_tier
    return {
        "type": "GOVERNANCE_ROOT_COMMIT",
        "governance_version": "v0.2.0",   # bumped at ADR-039
        "merkle_root": merkle_root,
        "drift_score": round(drift_score, 6),
        "drift_tier": classify_drift_tier(drift_score),
        "immutable_files": immutable_files,
        "signed_by": signed_by,
        "timestamp": _utcnow_iso(),
        "previous_root_tx_id": previous_tx_id,
        "governance_notes": notes,
    }
\end{lstlisting}

\subsection{Kill-Switch Binding (ADR-040)}

\begin{lstlisting}[
  style=pythonstyle,
  caption={Proposed \texttt{rollback/kill\_switch.py} drift trigger --- ADR-040 addition},
  label={lst:kill-switch}
]
# ADR-040: Bind kill_switch to PEET drift Collapse tier.
# This function is called by the CI drift monitor when
# classify_drift_tier() returns 'collapse'.
# HUMAN APPROVAL REQUIRED before execution (ADR-011 compliance).

from pirtm.constants import PEET_DRIFT_EPSILON
from pirtm.core.drift import drift_score, classify_drift_tier
from pirtm.governance.telemetry import telemetry


def drift_collapse_handler(
    n: int,
    psi: float,
    tensor_fn,
    t: float,
    dry_run: bool = True,
) -> dict:
    """
    ADR-040 Collapse tier handler.

    Evaluates current drift and, if Collapse, prepares kill-switch
    payload. Does NOT execute autonomously (dry_run=True default).
    A human custodian must call with dry_run=False after review.

    Returns:
        dict with keys: tier, delta, action, requires_approval
    """
    delta = drift_score(n, psi, tensor_fn, t)
    tier = classify_drift_tier(delta)

    result = {
        "tier": tier,
        "delta": delta,
        "epsilon": PEET_DRIFT_EPSILON,
        "action": None,
        "requires_approval": True,
        "dry_run": dry_run,
    }

    if tier == "collapse":
        result["action"] = "VELOCITY_STOP"
        telemetry.counter("pirtm_kill_switch_trigger_total",
                          labels={"dry_run": str(dry_run)})

        if not dry_run:
            # Custodian-approved execution path
            _execute_velocity_stop()
            result["executed"] = True
        else:
            result["executed"] = False
            result["message"] = (
                "Collapse tier reached. "
                "Human custodian approval required. "
                "Re-invoke with dry_run=False to execute."
            )

    return result


def _execute_velocity_stop() -> None:
    """Halt all active governance transitions. Irreversible."""
    import os, signal
    telemetry.counter("pirtm_velocity_stop_total")
    # Signal daemon to halt; implementation depends on daemon protocol
    os.kill(os.getpid(), signal.SIGTERM)
\end{lstlisting}

\subsection{Unit Tests (ADR-038 Exit Criterion)}

\begin{lstlisting}[
  style=pythonstyle,
  caption={\texttt{pirtm/tests/test\_drift.py} --- ADR-038 exit criterion tests},
  label={lst:tests}
]
"""ADR-038 exit criterion: drift_score() returns float for 3 known inputs."""

import pytest
from pirtm.core.drift import drift_score, classify_drift_tier, prime_factors


def identity_tensor(pi: int, pj: int, t: float) -> float:
    """Identity tensor: T(p_i, p_j, t) = 1 for all inputs."""
    return 1.0


def test_prime_factors_known():
    assert prime_factors(12) == [2, 2, 3]
    assert prime_factors(7)  == [7]
    assert prime_factors(30) == [2, 3, 5]


def test_drift_score_zero_case():
    # n=6, primes=[2,3], Psi = 2*2 + 2*3 + 3*2 + 3*3 = 4+6+6+9 = 25
    # psi = 25.0 -> delta = 0.0
    n, psi, t = 6, 25.0, 0.0
    delta = drift_score(n, psi, identity_tensor, t)
    assert isinstance(delta, float)
    assert delta == pytest.approx(0.0, abs=1e-9)


def test_drift_score_nonzero():
    # n=6, Psi=25.0, psi=28.0 -> delta=3.0
    n, psi, t = 6, 28.0, 0.0
    delta = drift_score(n, psi, identity_tensor, t)
    assert isinstance(delta, float)
    assert delta == pytest.approx(3.0, abs=1e-9)


def test_drift_score_prime_singleton():
    # n=7, primes=[7], Psi = 7*7 = 49
    n, psi, t = 7, 52.0, 1.0
    delta = drift_score(n, psi, identity_tensor, t)
    assert delta == pytest.approx(3.0, abs=1e-9)


def test_classify_nominal():
    assert classify_drift_tier(0.02) == "nominal"

def test_classify_watch():
    assert classify_drift_tier(0.07) == "watch"

def test_classify_warn():
    assert classify_drift_tier(0.12) == "warn"

def test_classify_collapse():
    # Assumes PEET_DRIFT_EPSILON < 0.35
    assert classify_drift_tier(0.35) == "collapse"


def test_drift_invalid_n():
    with pytest.raises(ValueError):
        drift_score(0, 1.0, identity_tensor, 0.0)


def test_kappa_out_of_tolerance():
    with pytest.raises(ValueError):
        drift_score(6, 25.0, identity_tensor, 0.0, kappa=1.5)
\end{lstlisting}

% ============================================================
\section{Prior Art Claims Register}
% ============================================================

This section constitutes the formal prior art disclosure registry.
Each claim is stated with sufficient specificity to establish
novelty and to prevent downstream patenting by any third party.

\subsection{Claim 1: PEET Drift Scalar as Governance Metric}

\textbf{Disclosed concept:} The use of Equation~\eqref{eq:peet-drift}
--- a prime-indexed entropy deviation scalar derived from the Prime
Entanglement Entropy Tensor --- as the canonical drift metric
governing software governance system state transitions, specifically
as the quantity recorded in a Merkle-chained governance ledger at
each commit boundary.

\textbf{Novelty basis:} No prior art known combining prime number
theory (prime factor tensor decomposition) with governance ledger
immutability chains for software system drift measurement.

\textbf{Disclosure date:} March 25, 2026 (initial ADR plan),
April 25, 2026 (this publication).

\subsection{Claim 2: Three-Tier Drift Response with Immutable Kill-Switch Binding}

\textbf{Disclosed concept:} A governance enforcement architecture
in which a continuously computed entropy drift scalar is classified
into three ordered response tiers (Watch, Warn, Collapse), where
the Collapse tier is bound to an immutable file registered in a
Merkle governance root, and execution requires human custodian
approval (precluding autonomous action per ADR-011).

\textbf{Novelty basis:} The combination of (a) entropy-based tier
classification, (b) immutable-file binding for the response artifact,
and (c) human-gated autonomous execution in a governance ledger
system constitutes a novel architectural pattern.

\subsection{Claim 3: Langlands Admissibility as Software State Gate}

\textbf{Disclosed concept:} The application of the
Langlands--modular form correspondence (\(Z_{\peet}(n) \in \modforms\))
as a formal admissibility predicate for software system state
transitions, implemented as a computable approximant in a spectral
module (\texttt{pirtm/spectral/}) that feeds into a runtime drift
measurement pipeline.

\textbf{Novelty basis:} No prior art known applying Langlands
duality or modular form membership tests as runtime software
governance gates.

\subsection{Claim 4: WAC / PEET Cohabitation in Configuration}

\textbf{Disclosed concept:} The co-residence of two incommensurable
drift proxies --- a Weighted Arithmetic Contractivity (WAC) ratio
(\texttt{maxWacAllowed}) and a PEET entropy norm threshold
(\texttt{PEET\_DRIFT\_EPSILON}) --- within the same configuration
artifact (\texttt{.phase-mirror.yml}), with explicit documentation
that they answer different questions and are not interchangeable,
as established by Theorem~\ref{thm:incommensurable}.

\subsection{Claim 5: Prime-Indexed Tensor Constants Module with Policy Loading}

\textbf{Disclosed concept:} A constants module (\texttt{pirtm/constants.py})
that loads L0 thresholds from a YAML policy artifact at import time,
exposes a two-layer threshold architecture (L0 operator contractivity
vs.\ L1 marshalling fidelity), and emits Prometheus-compatible
telemetry for each loaded constant, with explicit case analysis
of the four stability/fidelity combinations.

% ============================================================
\section{Recommended Future Work}
% ============================================================

\begin{enumerate}[leftmargin=2em]
  \item \textbf{Daemon real-time sampling.} Once ADR-039 establishes
        the batch-on-commit baseline, a subsequent ADR should bind
        \texttt{daemon/} to the \texttt{drift\_score()} API for
        sub-second telemetry, enabling the Constitutional Hamiltonian's
        3.33\,ms sampling cadence.

  \item \textbf{Coq/Idris formal proofs.} The Langlands admissibility
        gate (ADR-041) should ultimately be backed by a machine-checked
        proof that the computable approximant \(\hat{Z}_{\peet}(n)\)
        converges to the true modular form membership predicate within
        a certified error bound.

  \item \textbf{Metric unification.} A bridge function mapping
        \(\delta_H(t)\) (Frobenius norm) to an equivalent PEET
        representation should be specified, even if proven
        incommensurable in general, for the specific case of
        knot Hamiltonians with prime-indexed eigenstructure.

  \item \textbf{Quantum circuit PEET gates.} The \xicons{}
        Appendix D specifies a Qiskit PEET circuit. This should
        be implemented in \texttt{pirtm/} and registered as a
        \texttt{drift\_score()} backend for quantum runtime targets.

  \item \textbf{Exponential penalty enforcement.} The ethical drift
        penalty \(c(\tau) = c_0 e^{\lambda \tau}\) from the
        \xicons{} has no current runtime implementation. A penalty
        accumulator module should be designed that increments at
        the heartbeat cadence for states in the Collapse tier.
\end{enumerate}

% ============================================================
\section{Document Provenance and Integrity}
% ============================================================

\begin{center}
\begin{tabular}{ll}
\toprule
\textbf{Field} & \textbf{Value} \\
\midrule
Document ID           & \texttt{MF-DR-2026-001} \\
Repository            & \texttt{MultiplicityFoundation/PhaseMirror-HQ} \\
Branch (analysis)     & \texttt{main} @ \texttt{324a44bf...} \\
Governance ledger     & tx\#5, Merkle root \texttt{73e8aadf...} \\
Ledger bootstrap date & 2026-03-14T22:37:34Z \\
Last ledger rotation  & 2026-03-16T01:34:20Z \\
ADR high-water mark   & ADR-037 (at analysis date) \\
Proposed ADRs         & ADR-038, ADR-039, ADR-040, ADR-041 \\
License (code)        & Apache 2.0 \\
License (IP)          & Prime Materia License \\
Publication date      & April 25, 2026 \\
Author                & Ryan Van Gelder \\
Organization          & Multiplicity Foundation, 501(c)(3) \\
Location              & Bridgeville, Pennsylvania, USA \\
\bottomrule
\end{tabular}
\end{center}

\vspace{1em}
\begin{defensivebox}
\textbf{DEFENSIVE PUBLICATION ATTESTATION}\\[0.5em]
I, Ryan Van Gelder, hereby attest that the technical disclosures
contained in this document represent original work of the Multiplicity
Foundation, first conceived and implemented in the period
\textbf{March--April 2026}, and are intentionally published to
establish prior art. This document is submitted as a defensive
publication under the norms of open-source intellectual property
practice. No patent application covering the claims in Section~7
has been filed by the Multiplicity Foundation as of the publication
date, and none is intended. Third parties are hereby on notice that
these concepts constitute prior art.\\[0.5em]
\textbf{Signature:} \underline{\hspace{6cm}}
\quad \textbf{Date:} April 25, 2026
\end{defensivebox}

% ============================================================
%  APPENDIX
% ============================================================
\appendix

\section{Glossary}

\begin{description}[leftmargin=3cm, style=nextline]
  \item[ADR] Architecture Decision Record --- a document capturing
    a significant architectural decision, its context, and consequences.
  \item[PEET] Prime Entanglement Entropy Tensor --- the prime-indexed
    entropy field governing lawfulness of state transitions.
  \item[PIRTM] Prime Interval Recursion Theory Machine --- the core
    runtime of \pmhq{}.
  \item[WAC] Weighted Arithmetic Contractivity --- the discrete
    rule-violation ratio tracked by \texttt{maxWacAllowed}.
  \item[Drift Custodian] Named role (ADR-040) responsible for
    authorizing kill-switch execution on Collapse tier detection.
  \item[Collapse tier] The regime \(\delta_{\peet} > \varepsilon_{\peet}\)
    requiring human-gated kill-switch invocation.
  \item[Merkle root] A hash of all immutable governance files,
    forming a tamper-evident chain in \texttt{governance/ledger.json}.
  \item[Langlands gate] The admissibility predicate
    \(Z_{\peet}(n) \in \modforms\) derived from the Langlands program.
  \item[ModForms] The space of holomorphic modular forms over
    \(\mathrm{SL}_2(\mathbb{Z})\) or appropriate congruence subgroup.
  \item[L0 invariant] An operator-level mathematical stability
    condition that cannot be overridden by governance policy.
\end{description}

\section{ADR Dependency Graph}

\begin{center}
\begin{tabular}{ccccc}
ADR-038 & $\longrightarrow$ & ADR-039 & $\longrightarrow$ & ADR-040 \\
\multicolumn{5}{c}{$\downarrow$ \quad (all three required)} \\
\multicolumn{5}{c}{ADR-041} \\
\end{tabular}
\end{center}

\noindent Each arrow denotes a hard dependency: the downstream ADR
cannot enter \emph{Accepted} status until the upstream ADR exits
\emph{Draft} and its exit criterion is verified.

\section{Governance Ledger Hash Chain}

\begin{center}
\small
\begin{tabular}{clll}
\toprule
\textbf{Tx} & \textbf{Merkle Root (prefix)} & \textbf{Signed By} & \textbf{Timestamp} \\
\midrule
1 & \texttt{0958aeac1f29b760...} & bootstrap & 2026-03-14T22:37Z \\
2 & \texttt{54d0df56bd285e29...} & root-rotation-2026-03-16 & 2026-03-16T01:17Z \\
3 & \texttt{dbedae94e5676bd8...} & root-rotation-2026-03-16-d10 & 2026-03-16T01:24Z \\
4 & \texttt{73e8aadf105ee618...} & root-rotation-deterministic-2026-03-16 & 2026-03-16T01:30Z \\
5 & \texttt{73e8aadf105ee618...} & root-rotation-hash-semantics-2026-03-16 & 2026-03-16T01:34Z \\
\midrule
6 & \textit{(pending ADR-041)} & \textit{(TBD)} & \textit{(TBD)} \\
\bottomrule
\end{tabular}
\end{center}

\noindent Note: Transactions 4 and 5 share the same Merkle root
(\texttt{73e8aadf...}), indicating a hash-semantics canonicalization
rotation without file content change. This is consistent with the
\texttt{hash\_semantics: "canonicalized\_governance\_pointer"} field
present in tx\#5's \texttt{contracts/shared/constants.py} entry.

% ============================================================
%  MATHEMATICAL APPENDIX
%  MF-DR-2026-001 Amendment A
%  PEET Drift Protocol: Explicit Proofs and Operator Norm Bounds
%
%  Compile standalone:  pdflatex appendix-math-standalone.tex
%  Or \input into main: place after \appendix in MF-DR-2026-001.tex
%
%  Author:  Ryan Van Gelder, Multiplicity Foundation
%  Date:    April 25, 2026
%  Version: 1.0.0
% ============================================================

% ----------------------------------------------------------------
%  STANDALONE PREAMBLE (comment out if \input-ing into main doc)
% ----------------------------------------------------------------
\documentclass[11pt,letterpaper]{article}
\usepackage[margin=1in]{geometry}
\usepackage{amsmath,amssymb,amsthm,mathtools}
\usepackage{hyperref}
\usepackage{xcolor}
\usepackage{booktabs}
\usepackage{enumitem}
\usepackage{fancyhdr}
\usepackage{titlesec}
\usepackage{mdframed}
\usepackage{thmtools}
\usepackage{microtype}

\hypersetup{
  pdftitle={Mathematical Appendix: PEET Drift Protocol},
  pdfauthor={Ryan Van Gelder, Multiplicity Foundation},
  colorlinks=true, linkcolor=black, urlcolor=blue!70!black
}

\pagestyle{fancy}
\fancyhf{}
\rhead{\small MF-DR-2026-001 Amendment A}
\lhead{\small Mathematical Appendix: PEET Drift Protocol}
\cfoot{\thepage}
\renewcommand{\headrulewidth}{0.4pt}

% Theorem environments
\theoremstyle{plain}
\newtheorem{theorem}{Theorem}[section]
\newtheorem{lemma}[theorem]{Lemma}
\newtheorem{proposition}[theorem]{Proposition}
\newtheorem{corollary}[theorem]{Corollary}

\theoremstyle{definition}
\newtheorem{definition}[theorem]{Definition}
\newtheorem{invariant}[theorem]{Invariant}
\newtheorem{axiom}[theorem]{Axiom}
\newtheorem{construction}[theorem]{Construction}

\theoremstyle{remark}
\newtheorem{remark}[theorem]{Remark}
\newtheorem{example}[theorem]{Example}

% Custom operators
\DeclareMathOperator{\Spec}{Spec}
\DeclareMathOperator{\Tr}{Tr}
\DeclareMathOperator{\rank}{rank}
\DeclareMathOperator{\im}{im}
\DeclareMathOperator{\diam}{diam}
\newcommand{\peet}{\mathrm{PEET}}
\newcommand{\pirtm}{\textsc{PIRTM}}
\newcommand{\norm}[1]{\left\|#1\right\|}
\newcommand{\normF}[1]{\left\|#1\right\|_F}
\newcommand{\norminf}[1]{\left\|#1\right\|_\infty}
\newcommand{\abs}[1]{\left|#1\right|}
\newcommand{\inner}[2]{\langle #1,\, #2 \rangle}
\newcommand{\Psi}{\varPsi}
\newcommand{\eps}{\varepsilon}
\newcommand{\modforms}{\mathrm{ModForms}}
\newcommand{\specz}{\Spec\,\mathbb{Z}}
\newcommand{\R}{\mathbb{R}}
\newcommand{\N}{\mathbb{N}}
\newcommand{\Z}{\mathbb{Z}}
\newcommand{\C}{\mathbb{C}}
\newcommand{\Primes}{\mathbb{P}}
\newcommand{\calT}{\mathcal{T}}
\newcommand{\calM}{\mathcal{M}}
\newcommand{\calL}{\mathcal{L}}
\newcommand{\calH}{\mathcal{H}}
\newcommand{\calB}{\mathcal{B}}
\newcommand{\bbone}{\mathbf{1}}
\newcommand{\drift}{\delta_{\peet}}

% Boxed important results
\newmdenv[
  topline=true,bottomline=true,leftline=true,rightline=true,
  linewidth=1pt, linecolor=black!50,
  backgroundcolor=gray!4,
  innertopmargin=6pt,innerbottommargin=6pt,
  innerleftmargin=8pt,innerrightmargin=8pt,
  skipabove=8pt,skipbelow=8pt
]{importantbox}

\begin{document}

\begin{center}
  {\LARGE\bfseries Mathematical Appendix}\\[0.4em]
  {\large Explicit Proofs and Operator Norm Bounds}\\[0.3em]
  {\normalsize Amendment A to MF-DR-2026-001}\\[0.2em]
  {\small Ryan Van Gelder \quad Multiplicity Foundation \quad April 25, 2026}
\end{center}

\tableofcontents
\bigskip

\noindent\textbf{Scope.} This appendix provides complete proofs for all
theorems stated without proof in the main report, derives the operator
norm bounds used by the \texttt{validate\_operator\_contractivity()}
function (ADR-010), establishes the metric incommensurability result
(Theorem~\ref{thm:incommensurable-full}), and develops the spectral
radius bounds for the PEET drift pipeline. All constants cited
(\texttt{OPERATOR\_CONTRACTIVITY\_BOUND = 1.0},
\texttt{DEFAULT\_EPSILON = 0.05},
\texttt{DRIFT\_THRESHOLD = 0.3},
\texttt{MARSHALLING\_FIDELITY\_THRESHOLD = 0.95},
\texttt{SPECTRAL\_CONSERVATION\_TOLERANCE = 0.02}) are sourced
directly from \texttt{packages/pirtm/constants.py} (ADR-010).

% ============================================================
\section{Algebraic Foundations of Multiplicity Spaces}
\label{sec:foundations}
% ============================================================

\subsection{Prime Decomposition Embedding}

\begin{definition}[Prime Basis]
\label{def:prime-basis}
Let $\Primes = (p_1, p_2, p_3, \ldots) = (2, 3, 5, 7, 11, \ldots)$
be the ordered sequence of primes. For any $n \in \N_{>0}$, write
its canonical prime factorization as
\[
  n = \prod_{k \geq 1} p_k^{\alpha_k(n)},
  \quad \alpha_k(n) \in \N_{\geq 0}, \quad \text{finitely many nonzero.}
\]
The \emph{prime signature} of $n$ is the vector
$\boldsymbol{\alpha}(n) = (\alpha_1(n), \alpha_2(n), \ldots) \in \N^{\infty}$.
\end{definition}

\begin{definition}[Prime Decomposition Embedding]
\label{def:phi}
The \emph{prime decomposition embedding} is the map
$\phi: \N_{>0} \to \ell^2(\Primes)$ defined by
\[
  \phi(n) = \sum_{p \mid n} \alpha_p(n) \cdot e_p
\]
where $e_p$ is the standard basis vector at position $p$, and
$\alpha_p(n) = v_p(n)$ is the $p$-adic valuation of $n$.
\end{definition}

\begin{proposition}[Isometric Embedding into $\ell^2$]
\label{prop:isometric}
The map $\phi$ satisfies:
\begin{enumerate}[label=(\roman*),noitemsep]
  \item $\phi(mn) = \phi(m) + \phi(n)$ for all $m,n \in \N_{>0}$
        (i.e., $\phi$ is a monoid homomorphism $(\N_{>0}, \times) \to (\ell^2, +)$).
  \item $\norm{\phi(n)}_{\ell^2}^2 = \sum_{p \mid n} v_p(n)^2 = \Omega_2(n)$,
        the sum of squares of prime exponents of $n$.
  \item $\phi$ is injective: $\phi(m) = \phi(n) \Rightarrow m = n$
        (follows from the Fundamental Theorem of Arithmetic).
\end{enumerate}
\end{proposition}

\begin{proof}
(i) By definition,
$\phi(mn)_p = v_p(mn) = v_p(m) + v_p(n) = \phi(m)_p + \phi(n)_p$
for all $p \in \Primes$.
(ii) Immediate from $\norm{\phi(n)}_{\ell^2}^2 = \sum_p v_p(n)^2$.
(iii) If $\phi(m) = \phi(n)$, then $v_p(m) = v_p(n)$ for all $p$,
so $m = n$ by unique factorization. \qed
\end{proof}

% ============================================================
\section{Prime-Interaction Tensor: Construction and Bounds}
\label{sec:tensor}
% ============================================================

\subsection{Tensor Definition and Regularity Conditions}

\begin{definition}[Prime-Interaction Tensor]
\label{def:tensor}
A \emph{prime-interaction tensor} at time $t \geq 0$ is a
function $\calT: \Primes \times \Primes \times \R_{\geq 0} \to \R$
satisfying:
\begin{enumerate}[label=(T\arabic*),noitemsep]
  \item\label{T1} \textbf{Symmetry:} $\calT(p,q,t) = \calT(q,p,t)$
        for all $p,q,t$.
  \item\label{T2} \textbf{Boundedness:} $\sup_{p,q,t} \abs{\calT(p,q,t)} \leq M_{\calT} < \infty$.
  \item\label{T3} \textbf{Lipschitz continuity in $t$:}
        $\abs{\calT(p,q,t) - \calT(p,q,s)} \leq L_{\calT} \abs{t-s}$
        for all $p,q$ and all $t,s \geq 0$.
  \item\label{T4} \textbf{Summability:} For each fixed $t$,
        $\sum_{p,q \in P(n)} \abs{\calT(p,q,t)} \cdot pq < \infty$
        for all $n \in \N_{>0}$.
\end{enumerate}
\end{definition}

\begin{lemma}[Tensor Norm Bound]
\label{lem:tensor-norm}
Let $\calT$ satisfy conditions \ref{T1}--\ref{T4}, and let
$P(n) = \{p_1, \ldots, p_k\}$ be the distinct prime factors of $n$
(with $k = \omega(n)$ the number of distinct prime divisors).
Then the prime-tensor superposition satisfies:
\[
  \abs{\Psi(n,t)} \leq M_{\calT} \cdot \left(\sum_{p \in P(n)} p\right)^2
  \leq M_{\calT} \cdot \omega(n)^2 \cdot p_{\max}(n)^2
\]
where $p_{\max}(n)$ is the largest prime factor of $n$.
\end{lemma}

\begin{proof}
By definition,
\[
  \abs{\Psi(n,t)}
  = \abs[\bigg]{\sum_{p_i, p_j \in P(n)} \calT(p_i,p_j,t) \cdot p_i p_j}
  \leq \sum_{p_i, p_j \in P(n)} \abs{\calT(p_i,p_j,t)} \cdot p_i p_j.
\]
Applying condition \ref{T2}:
\[
  \leq M_{\calT} \sum_{p_i \in P(n)} \sum_{p_j \in P(n)} p_i p_j
  = M_{\calT} \left(\sum_{p \in P(n)} p\right)^2.
\]
The second inequality uses $\sum_{p \in P(n)} p \leq \omega(n) \cdot p_{\max}(n)$. \qed
\end{proof}

\subsection{Time-Evolution Stability of the Tensor}

\begin{lemma}[Temporal Drift of Superposition]
\label{lem:temporal-drift}
Under conditions \ref{T1}--\ref{T4}, the superposition $\Psi(n,\cdot)$
is Lipschitz continuous in $t$ with constant:
\[
  \abs{\Psi(n,t) - \Psi(n,s)} \leq L_{\calT} \cdot \omega(n)^2 \cdot p_{\max}(n)^2 \cdot \abs{t-s}.
\]
\end{lemma}

\begin{proof}
\begin{align*}
\abs{\Psi(n,t) - \Psi(n,s)}
&= \abs[\bigg]{\sum_{p_i,p_j \in P(n)}
    \bigl(\calT(p_i,p_j,t) - \calT(p_i,p_j,s)\bigr) p_i p_j} \\
&\leq \sum_{p_i,p_j \in P(n)}
    \abs{\calT(p_i,p_j,t) - \calT(p_i,p_j,s)} \cdot p_i p_j \\
&\leq L_{\calT} \abs{t-s} \sum_{p_i,p_j \in P(n)} p_i p_j
    \quad \text{(condition \ref{T3})} \\
&= L_{\calT} \abs{t-s} \left(\sum_{p \in P(n)} p\right)^2
\leq L_{\calT} \cdot \omega(n)^2 \cdot p_{\max}(n)^2 \cdot \abs{t-s}. \qed
\end{align*}
\end{proof}

% ============================================================
\section{The PEET Drift Scalar: Full Analysis}
\label{sec:drift-analysis}
% ============================================================

\subsection{Existence and Continuity}

\begin{theorem}[Well-Posedness of Drift Scalar]
\label{thm:well-posed}
Let $\psi: \N_{>0} \times \R_{\geq 0} \to \R$ be a state function
satisfying $\abs{\psi(n,t)} \leq C_\psi$ for all $n,t$, and let
$\calT$ satisfy \ref{T1}--\ref{T4}. Then for every $n \in \N_{>0}$
and $t \geq 0$:
\begin{enumerate}[label=(\roman*),noitemsep]
  \item The drift scalar
        $\drift(n,t) = \abs{\psi(n,t) - \Psi(n,t)}$
        is well-defined and finite.
  \item $\drift(n,\cdot)$ is Lipschitz continuous in $t$ with constant:
        \[
          \Lip_t(\drift) \leq L_\psi + L_{\calT}\,\omega(n)^2\,p_{\max}(n)^2
        \]
        where $L_\psi$ is the Lipschitz constant of $\psi(n,\cdot)$.
  \item The zero-drift condition $\drift(n,t) = 0$ holds if and only if
        $\psi(n,t) = \Psi(n,t)$, i.e., the wavefunction equals
        its prime-tensor superposition exactly.
\end{enumerate}
\end{theorem}

\begin{proof}
(i) Finiteness: $\abs{\psi(n,t)} \leq C_\psi < \infty$ by assumption,
and $\abs{\Psi(n,t)} \leq M_{\calT} \cdot \omega(n)^2 \cdot p_{\max}(n)^2 < \infty$
by Lemma~\ref{lem:tensor-norm}. Hence
$\drift(n,t) \leq C_\psi + M_{\calT}\,\omega(n)^2\,p_{\max}(n)^2 < \infty$.

(ii) For any $t, s \geq 0$, apply the reverse triangle inequality:
\begin{align*}
  \abs{\drift(n,t) - \drift(n,s)}
  &= \abs[\big]{\abs{\psi(n,t)-\Psi(n,t)} - \abs{\psi(n,s)-\Psi(n,s)}} \\
  &\leq \abs{(\psi(n,t) - \Psi(n,t)) - (\psi(n,s) - \Psi(n,s))} \\
  &\leq \abs{\psi(n,t) - \psi(n,s)} + \abs{\Psi(n,t) - \Psi(n,s)} \\
  &\leq L_\psi \abs{t-s} + L_{\calT}\,\omega(n)^2\,p_{\max}(n)^2\,\abs{t-s}
\end{align*}
using Lemma~\ref{lem:temporal-drift} for the second term.

(iii) Immediate from $\drift(n,t) = \abs{\psi(n,t) - \Psi(n,t)} = 0
\Leftrightarrow \psi(n,t) = \Psi(n,t)$. \qed
\end{proof}

\subsection{Drift Under Spectral Curvature Correction}

\begin{theorem}[Spectral Correction Bound]
\label{thm:spectral-kappa}
Let $\kappa \in [1 - \sigma, 1 + \sigma]$ where
$\sigma = \texttt{SPECTRAL\_CONSERVATION\_TOLERANCE} = 0.02$.
The corrected drift scalar
$\drift^{(\kappa)}(n,t) = \abs{\psi(n,t) - \Psi(n,t)} \cdot \kappa$
satisfies:
\[
  (1-\sigma)\,\drift(n,t)
  \;\leq\;
  \drift^{(\kappa)}(n,t)
  \;\leq\;
  (1+\sigma)\,\drift(n,t).
\]
In particular, the tier classification from Section~3.6 of the main
report is stable under $\kappa$-correction provided the tier thresholds
carry a safety margin of at least $\sigma \cdot \eps_X$ at each boundary.
\end{theorem}

\begin{proof}
Since $\kappa \in [1-\sigma, 1+\sigma]$ and $\drift(n,t) \geq 0$:
\[
  (1-\sigma)\,\drift(n,t)
  = \drift(n,t) \cdot (1-\sigma)
  \leq \drift(n,t) \cdot \kappa
  = \drift^{(\kappa)}(n,t)
  \leq \drift(n,t) \cdot (1+\sigma).
\]
For tier stability: a state with $\drift(n,t) \leq \eps_X - \sigma\eps_X
= (1-\sigma)\eps_X$ satisfies $\drift^{(\kappa)}(n,t) \leq (1+\sigma)
\cdot (1-\sigma)\eps_X = (1-\sigma^2)\eps_X < \eps_X$ for all $\kappa$
in the tolerance band. Hence tier boundaries are preserved under small
spectral correction. \qed
\end{proof}

% ============================================================
\section{Operator Norm Bounds and L0 Contractivity}
\label{sec:operator-norms}
% ============================================================

This section provides the rigorous derivation underlying
\texttt{validate\_operator\_contractivity()} in ADR-010.

\subsection{The Composite Evolution Operator}

\begin{definition}[PIRTM Evolution Operator]
\label{def:evolution-op}
Let $\calH$ be a separable Hilbert space of system states.
The \emph{PIRTM composite evolution operator} at time $t$ is:
\[
  \Xi_t : \calH \to \calH,
  \qquad
  \Xi_t = \Xi + \Lambda \circ \mathcal{F}_T
\]
where:
\begin{itemize}[noitemsep,leftmargin=2em]
  \item $\Xi: \calH \to \calH$ is the base state-transition operator,
  \item $\Lambda: \calH \to \calH$ is the Lipschitz feedback operator,
  \item $\mathcal{F}_T: \calH \to \calH$ is the prime-tensor folding map
        associated to tensor $\calT$ at time $t$,
  \item $L_T \defeq \Lip(\mathcal{F}_T)$ is the Lipschitz constant of the
        tensor folding (from \texttt{validate\_operator\_contractivity}
        docstring: \texttt{q\_t = \|\Xi\| + \|\Lambda\| * L\_T}).
\end{itemize}
\end{definition}

\begin{definition}[Composite Spectral Radius]
The \emph{composite spectral radius} $q_t$ is:
\begin{equation}
  q_t \;\defeq\; \norm{\Xi}_{\calB(\calH)} + \norm{\Lambda}_{\calB(\calH)} \cdot L_T
  \label{eq:qt}
\end{equation}
where $\norm{\cdot}_{\calB(\calH)}$ denotes the operator norm on
bounded linear operators $\calH \to \calH$.
\end{definition}

\subsection{The L0 Contractivity Theorem}

\begin{theorem}[L0 Contractivity --- Formal Statement]
\label{thm:l0-contractivity}
Let $\Xi_t$ be as in Definition~\ref{def:evolution-op} with composite
spectral radius $q_t$ as in~\eqref{eq:qt}. Let
$B = \texttt{OPERATOR\_CONTRACTIVITY\_BOUND} = 1.0$ and
$\eps = \texttt{DEFAULT\_EPSILON} = 0.05$.
Define the \emph{effective contractivity bound}
$B_\eps \defeq B - \eps = 0.95$.

If $q_t < B_\eps$, then $\Xi_t$ is a \emph{strict contraction} on $\calH$:
\begin{enumerate}[label=(\roman*),noitemsep]
  \item $\norm{\Xi_t x - \Xi_t y}_\calH \leq q_t \norm{x-y}_\calH$
        for all $x, y \in \calH$.
  \item There exists a unique fixed point $x^* \in \calH$ with
        $\Xi_t(x^*) = x^*$.
  \item The fixed point iteration $x_{k+1} = \Xi_t(x_k)$ converges
        geometrically:
        \[
          \norm{x_k - x^*}_\calH \leq \frac{q_t^k}{1-q_t}\,\norm{x_1 - x_0}_\calH.
        \]
\end{enumerate}
\end{theorem}

\begin{proof}
The operator $\Xi_t = \Xi + \Lambda \circ \mathcal{F}_T$ is a
composition of bounded linear (or Lipschitz) maps. For any $x, y \in \calH$:
\begin{align*}
  \norm{\Xi_t x - \Xi_t y}_\calH
  &= \norm{(\Xi x + \Lambda\mathcal{F}_T x) - (\Xi y + \Lambda\mathcal{F}_T y)}_\calH \\
  &\leq \norm{\Xi(x-y)}_\calH + \norm{\Lambda(\mathcal{F}_T x - \mathcal{F}_T y)}_\calH \\
  &\leq \norm{\Xi} \norm{x-y} + \norm{\Lambda} \cdot L_T \norm{x-y} \\
  &= q_t \norm{x-y}_\calH.
\end{align*}
Since $q_t < B_\eps < 1$, $\Xi_t$ is a strict contraction with
Lipschitz constant $q_t < 1$.
Parts (ii) and (iii) follow directly from the Banach Fixed-Point
(Contraction Mapping) Theorem, which applies since $(\calH, \norm{\cdot})$
is a complete metric space and $\Xi_t$ is a strict contraction. \qed
\end{proof}

\begin{corollary}[Instability When $q_t \geq B_\eps$]
\label{cor:instability}
If $q_t \geq B_\eps = 0.95$, the contraction guarantee fails.
In particular, if $q_t \geq B = 1.0$, the operator may be
\emph{expansive}: there exist $x \neq y$ such that
$\norm{\Xi_t x - \Xi_t y} \geq \norm{x-y}$, and no fixed-point
convergence is guaranteed.
\end{corollary}

\begin{remark}
The margin $\eps = 0.05$ (the gap between $B_\eps = 0.95$ and the
expansion threshold $B = 1.0$) provides a safety buffer against
perturbations in $\norm{\Xi}$ and $\norm{\Lambda} \cdot L_T$ due to
floating-point accumulation, tensor approximation error, or
governance ledger rotation events. This margin is consistent with
the 5\% marshalling fidelity threshold (\texttt{MARSHALLING\_FIDELITY\_THRESHOLD = 0.95})
at the L1 layer.
\end{remark}

\subsection{Operator Norm Decomposition}

\begin{proposition}[Submultiplicativity of $q_t$]
\label{prop:submult}
For a composed evolution $\Xi_t^{(N)} = \Xi_t \circ \Xi_t \circ \cdots \circ \Xi_t$
($N$ applications), the spectral radius satisfies:
\[
  \norm{\Xi_t^{(N)}}_{\calB(\calH)} \leq q_t^N \xrightarrow{N\to\infty} 0
  \quad \text{whenever } q_t < 1.
\]
\end{proposition}

\begin{proof}
By submultiplicativity of the operator norm:
$\norm{\Xi_t^{(N)}} \leq \norm{\Xi_t}^N \leq q_t^N$,
where the second inequality holds because
$\norm{\Xi_t} = \norm{\Xi + \Lambda\mathcal{F}_T} \leq
\norm{\Xi} + \norm{\Lambda} L_T = q_t$ by the triangle inequality. \qed
\end{proof}

\subsection{Perturbation Bound for Governance Rotations}

When the governance root is rotated (a new Merkle root is signed),
the operator may be perturbed by $\Delta\Xi_t$. The following
proposition quantifies the stability of the contractivity condition
under such perturbations.

\begin{proposition}[Stability Under Operator Perturbation]
\label{prop:perturbation}
Let $\Xi_t' = \Xi_t + \Delta$ be a perturbed operator with
$\norm{\Delta} \leq \delta_\Delta$. Then:
\[
  q_t' \defeq \norm{\Xi_t' + \Lambda\mathcal{F}_T}
  \leq q_t + \delta_\Delta.
\]
The perturbed system remains L0-contractive provided:
\[
  \delta_\Delta < B_\eps - q_t = 0.95 - q_t.
\]
\end{proposition}

\begin{proof}
$\norm{\Xi_t' + \Lambda\mathcal{F}_T} \leq \norm{\Xi_t + \Lambda\mathcal{F}_T} + \norm{\Delta}
= q_t + \delta_\Delta$.
The contractivity condition $q_t' < B_\eps$ then requires
$q_t + \delta_\Delta < 0.95$, i.e., $\delta_\Delta < 0.95 - q_t$. \qed
\end{proof}

\begin{remark}
This bound is directly relevant to governance ledger rotation events
in \pmhq{}: when a new Merkle root is signed (transactions \#1--\#5
in the ledger), any change to immutable files can perturb $\Xi_t$
by at most $\delta_\Delta = \norm{\Delta}_F$ (Frobenius norm of the
diff). Proposition~\ref{prop:perturbation} provides the certification
condition that must hold for the post-rotation system to remain L0-stable.
\end{remark}

% ============================================================
\section{Metric Incommensurability: Full Proof}
\label{sec:incommensurability}
% ============================================================

\begin{theorem}[Metric Incommensurability --- Complete Proof]
\label{thm:incommensurable-full}
Let $\drift(n,t)$ be the PEET drift scalar (Definition in main report,
eq.~(1)) and $\delta_H(t)$ be the Hamiltonian drift ratio
(main report eq.~(2)). There exists no $\alpha \in \R_{>0}$ and
no function $f: \N_{>0} \to \R_{>0}$ such that
\[
  \drift(n,t) = f(n) \cdot \delta_H(t)
\]
holds for all $n \in \N_{>0}$ and all $t \geq 0$.
\end{theorem}

\begin{proof}
We establish incommensurability on three independent axes.

\medskip
\noindent\textbf{(A) Domain structure.}
$\drift(n,t)$ depends on the prime-factor multiset $P(n)$, which is
a discrete combinatorial invariant of $n$. Specifically, for any
fixed $t$, the map $n \mapsto \drift(n,t)$ is not monotone:
e.g., for $n = 4 = 2^2$ vs.\ $n = 6 = 2 \cdot 3$, the
prime structure changes qualitatively (one prime vs.\ two primes).

By contrast, $\delta_H(t)$ is independent of $n$ entirely ---
it is a property of the Hamiltonian operator trajectory in an
abstract Hilbert space. Hence for any fixed $t$, the ratio
$\drift(n,t) / \delta_H(t)$ depends on $n$, meaning no
$n$-independent $\alpha$ can equate them.

\medskip
\noindent\textbf{(B) Unit analysis.}
The PEET drift scalar $\drift(n,t) = \abs{\psi(n,t) - \Psi(n,t)}$
has units of $[\psi]$ --- the same units as the state amplitude.
When $\psi$ is normalized to $[0,1]$, $\drift(n,t) \in [0,1]$ is
dimensionless in this normalization, but its scaling law with
respect to $n$ is governed by $\Psi(n,t) \sim M_{\calT}
\cdot (\sum_{p \in P(n)} p)^2$ (Lemma~\ref{lem:tensor-norm}),
which grows without bound as $n \to \infty$ through primes.

The Hamiltonian ratio $\delta_H(t) = \normF{H(t) - H^{(0)}} / E_\text{free}$
is bounded in $[0, \infty)$ independently of any integer argument.
Its scaling with respect to system parameters is determined by
Hilbert-space operator norm growth, not by prime factorization.

Therefore: $\drift(n,t) \to \infty$ as $n$ grows through numbers
with many large prime factors (e.g., $n = p_1 p_2 \cdots p_k$
for the first $k$ primes), while $\delta_H(t)$ is finite and
fixed for each $t$. No constant $\alpha$ reconciles this.

\medskip
\noindent\textbf{(C) Structural independence.}
Suppose for contradiction that $\drift(n,t) = f(n) \cdot \delta_H(t)$
for some $f: \N_{>0} \to \R_{>0}$. Fix $n$ and vary $t$.
The left side $\drift(n,t) = \abs{\psi(n,t) - \Psi(n,t)}$ is
Lipschitz in $t$ with a constant depending on $n$
(Theorem~\ref{thm:well-posed}(ii)). The right side $f(n) \cdot \delta_H(t)$
is Lipschitz in $t$ with a constant proportional to
$f(n) \cdot \normF{\dot{H}(t)}/E_\text{free}$.
These Lipschitz constants are determined by entirely different
dynamics (prime-tensor Lipschitz constant $L_{\calT}$ vs.\ Hamiltonian
flow rate $\normF{\dot{H}}$), and there is no physical or mathematical
principle coupling them. Choosing a concrete $\psi$ and $H$ that
are dynamically independent (e.g., $\psi$ driven by prime-sieve
updates and $H$ driven by knot-invariant evolution) yields
$\partial_t \drift(n,t) \neq f(n) \cdot \partial_t \delta_H(t)$
for generic $(n,t)$, contradicting the assumed equality. \qed
\end{proof}

\begin{corollary}[Three-Way Incommensurability]
\label{cor:three-way}
The \texttt{maxWacAllowed} ratio $\rho_{\wac} = \texttt{WAC}/\texttt{WAC}_0$
is a third incommensurable quantity. It is:
\begin{itemize}[noitemsep,leftmargin=2em]
  \item Dimensionless (ratio of two discrete rule-violation counts),
  \item Defined over a finite alphabet of MD-rule outcomes (not $\N_{>0}$),
  \item Bounded in $[0, \texttt{maxWacAllowed}] = [0, 1.05]$,
        a compact interval independent of prime structure or
        Hamiltonian dynamics.
\end{itemize}
No injection $\drift(n,t) \hookrightarrow \rho_{\wac}$ or
$\delta_H(t) \hookrightarrow \rho_{\wac}$ preserving order structure
can exist, since the domain of $\rho_{\wac}$ is finite while
both $\drift(n,t)$ and $\delta_H(t)$ are defined over infinite
parameter spaces.
\end{corollary}

% ============================================================
\section{Marshalling Fidelity and L1 Layer Bounds}
\label{sec:l1-bounds}
% ============================================================

\subsection{Round-Trip Fidelity}

\begin{definition}[Round-Trip Error]
For a marshalling path $\mu: X \to Y$ with inverse
$\mu^{-1}: Y \to X$, the \emph{round-trip error} for element
$x \in X$ with norm $\norm{x}_X > 0$ is:
\[
  e_\mu(x) = \frac{\norm{\mu^{-1}(\mu(x)) - x}_X}{\norm{x}_X}
  \in [0, \infty).
\]
The \emph{marshalling fidelity score} is:
\[
  s_\mu(x) = 1 - e_\mu(x) \in (-\infty, 1].
\]
\end{definition}

\begin{theorem}[L1 Fidelity Bound]
\label{thm:l1-fidelity}
A marshalling path $\mu$ is \emph{L1-acceptable} (passes the
\texttt{MARSHALLING\_FIDELITY\_THRESHOLD = 0.95} gate) if and only if:
\[
  s_\mu(x) \geq 0.95 \;\Longleftrightarrow\; e_\mu(x) \leq 0.05
  \;\Longleftrightarrow\; \norm{\mu^{-1}(\mu(x)) - x}_X \leq 0.05\,\norm{x}_X.
\]
This is equivalent to: the round-trip composition
$\mu^{-1} \circ \mu$ acts as the identity on $x$ up to a
relative perturbation of at most $5\%$.
\end{theorem}

\begin{proof}
Direct algebraic manipulation:
$s_\mu(x) \geq 0.95 \Leftrightarrow 1 - e_\mu(x) \geq 0.95
\Leftrightarrow e_\mu(x) \leq 0.05
\Leftrightarrow \norm{\mu^{-1}(\mu(x)) - x}/\norm{x} \leq 0.05
\Leftrightarrow \norm{\mu^{-1}(\mu(x)) - x} \leq 0.05\norm{x}$. \qed
\end{proof}

\subsection{Independence of L0 and L1 Conditions}

\begin{theorem}[L0/L1 Logical Independence]
\label{thm:l0l1-independence}
The L0 contractivity condition ($q_t < 0.95$) and the L1 fidelity
condition ($s_\mu(x) \geq 0.95$) are logically independent:
each can hold or fail independently of the other.
\end{theorem}

\begin{proof}
We exhibit four configurations (matching the case analysis in ADR-010):

\begin{enumerate}[label=\textbf{Case \arabic*:},leftmargin=3.5em]
  \item \textbf{L0 stable, L1 acceptable.}
        Take $\Xi = 0.9\,I$ (contraction), $\mu = I$ (identity, perfect fidelity).
        Then $q_t = 0.9 < 0.95$ and $s_\mu = 1.0 \geq 0.95$. Both pass.

  \item \textbf{L0 stable, L1 unacceptable.}
        Take $\Xi = 0.9\,I$ and $\mu$ such that $e_\mu = 0.08 > 0.05$.
        Then $q_t < 0.95$ (L0 passes) but $s_\mu = 0.92 < 0.95$ (L1 fails).

  \item \textbf{L0 unstable, L1 acceptable.}
        Take $\Xi = 1.1\,I$ (expansion), $\mu = I$.
        Then $q_t \geq 1.1 > 0.95$ (L0 fails) but $s_\mu = 1.0$ (L1 passes).
        This is the critical case: good communication cannot rescue
        an unstable operator.

  \item \textbf{L0 unstable, L1 unacceptable.}
        Take $\Xi = 1.1\,I$ and $e_\mu = 0.08$.
        Both conditions fail simultaneously.
\end{enumerate}

Since we have exhibited realizations of all four Boolean combinations,
the two conditions are logically independent. \qed
\end{proof}

% ============================================================
\section{Spectral Radius Conservation in Tensor Transforms}
\label{sec:spectral-conservation}
% ============================================================

\begin{definition}[Pi-to-Tensor Transform]
\label{def:pi-tensor}
A \emph{pi-to-tensor transform} is a linear map
$\Phi: \ell^2(\Primes) \to \calH_T$ from the prime-index space
to the tensor Hilbert space. The \emph{spectral conservation}
condition requires that the principal eigenvalue $\lambda_1(\Phi)$
satisfies:
\[
  \lambda_1(\Phi) \in
  [1 - \sigma,\; 1 + \sigma], \quad
  \sigma = \texttt{SPECTRAL\_CONSERVATION\_TOLERANCE} = 0.02.
\]
\end{definition}

\begin{theorem}[Spectral Conservation Bound]
\label{thm:spectral-conservation}
Let $\Phi$ be a pi-to-tensor transform with
$\lambda_1(\Phi) \in [1-\sigma, 1+\sigma]$.
Then for any $v \in \ell^2(\Primes)$ with $\norm{v} = 1$:
\[
  \abs{\norm{\Phi v}^2 - 1} \leq 2\sigma + \sigma^2 < 2\sigma + 1
\]
and in the regime $\sigma \ll 1$ (here $\sigma = 0.02$):
\[
  \abs{\norm{\Phi v}^2 - 1} \lesssim 2\sigma = 0.04.
\]
\end{theorem}

\begin{proof}
Let $\lambda = \lambda_1(\Phi) \in [1-\sigma, 1+\sigma]$.
For a unit eigenvector $v$: $\norm{\Phi v}^2 = \lambda^2$.
\[
  \abs{\lambda^2 - 1} = \abs{(\lambda-1)(\lambda+1)}
  \leq \sigma \cdot (2 + \sigma) = 2\sigma + \sigma^2.
\]
For $\sigma = 0.02$: $2\sigma + \sigma^2 = 0.04 + 0.0004 = 0.0404$.
For general unit vectors $v$ (not necessarily eigenvectors), the bound
holds for the dominant eigenvalue contribution; the full spectral
bound follows from the interlacing theorem applied to the
self-adjoint part of $\Phi^*\Phi$. \qed
\end{proof}

\begin{corollary}[Consistency of $\sigma$ and $\eps$]
\label{cor:sigma-eps-consistency}
The choices $\sigma = 0.02$ and $\eps = 0.05$ are mutually consistent
in the following sense: a spectral perturbation of magnitude $\sigma$
at the pi-to-tensor level induces an operator-norm perturbation of at
most $\delta_\Delta \leq 2\sigma = 0.04 < \eps = 0.05$. By
Proposition~\ref{prop:perturbation}, L0 contractivity is preserved
under pi-to-tensor transforms provided $q_t < B_\eps - 2\sigma = 0.93$.
\end{corollary}

% ============================================================
\section{PEET Partition Function and Modular Form Approximant}
\label{sec:langlands}
% ============================================================

\subsection{Formal q-Series Construction}

\begin{definition}[PEET Partition Function]
\label{def:z-peet}
For $n \in \N_{>0}$ and $q = e^{2\pi i \tau}$ with $\tau \in \calH$
(upper half-plane), define:
\[
  Z_{\peet}(n;\tau) = \sum_{k=1}^{\infty} a_k(n)\, q^k
\]
where the Fourier coefficients are:
\[
  a_k(n) = \sum_{\substack{p_i, p_j \in P(n) \\ p_i p_j = k}}
  \calT(p_i, p_j, 0) \cdot v_{p_i}(n) \cdot v_{p_j}(n).
\]
The sum ranges over all representations of $k$ as a product of
two prime factors of $n$ (with multiplicity from the valuations).
\end{definition}

\begin{remark}
For $k$ not expressible as $p_i p_j$ with $p_i, p_j \in P(n)$,
we have $a_k(n) = 0$. The series therefore has finite support at
each level determined by the prime structure of $n$.
\end{remark}

\subsection{Computable Approximant for ADR-041}

\begin{construction}[Finite-Level Approximant]
\label{const:approx}
For computational purposes (ADR-041), truncate the series at level
$K = p_{\max}(n)^2$:
\[
  \hat{Z}_{\peet}^{(K)}(n;\tau)
  = \sum_{k=1}^{K} a_k(n)\, q^k.
\]
This is a polynomial in $q$ of degree $K$, computable in $O(K \cdot \omega(n)^2)$
arithmetic operations. The approximation error is:
\[
  \abs{Z_{\peet}(n;\tau) - \hat{Z}_{\peet}^{(K)}(n;\tau)}
  \leq \sum_{k > K} \abs{a_k(n)} \cdot \abs{q}^k
  \leq \frac{M_{\calT} \cdot \omega(n)^2 \cdot \abs{q}^{K+1}}{1 - \abs{q}}
\]
for $\abs{q} < 1$, which holds since $\tau \in \calH$ implies
$\abs{q} = e^{-2\pi \Im(\tau)} < 1$.
\end{construction}

\begin{proposition}[Modular Form Proximity Criterion]
\label{prop:modform-criterion}
Let $f \in \modforms_w(\Gamma_0(N))$ be a weight-$w$ modular form
for the congruence subgroup $\Gamma_0(N)$ with Fourier expansion
$f(\tau) = \sum_{k\geq 1} b_k q^k$. The \emph{Langlands proximity
criterion} for $Z_{\peet}(n;\tau)$ is:
\[
  \mathcal{D}(n, f) \defeq
  \sum_{k=1}^{K} \frac{\abs{a_k(n) - b_k}^2}{1 + \abs{b_k}^2}
  < \theta_{\mathrm{Langlands}}
\]
for a threshold $\theta_{\mathrm{Langlands}} > 0$ to be calibrated
in ADR-041. When this criterion is satisfied, the state $n$ is
considered \emph{approximately Langlands-admissible} at resolution $K$.
\end{proposition}

% ============================================================
\section{Exponential Penalty: Formal Structure}
\label{sec:penalty}
% ============================================================

\begin{definition}[Ethical Drift Penalty Functional]
\label{def:penalty}
Let $\tau_C(t_0) = \inf\{t \geq t_0 : \drift(n,t) \leq \eps_C\}$
be the first exit time from the Collapse tier after time $t_0$.
The \emph{accumulated ethical penalty} for a trajectory remaining
in the Collapse tier from time $t_0$ to $t$ is:
\[
  C(t; t_0) = c_0 \int_{t_0}^{\min(t,\, \tau_C(t_0))}
  e^{\lambda(s - t_0)}\, ds
  = \frac{c_0}{\lambda}\left(e^{\lambda(t-t_0)} - 1\right)
\]
for constants $c_0 > 0$ (base penalty rate) and $\lambda > 0$
(growth exponent), valid while the system remains in the Collapse tier.
\end{definition}

\begin{proposition}[Penalty Growth Rate]
\label{prop:penalty-growth}
The penalty satisfies:
\begin{enumerate}[label=(\roman*),noitemsep]
  \item $C(t;t_0) = 0$ at $t = t_0$ (no instantaneous penalty).
  \item $\frac{d}{dt}C(t;t_0) = c_0 e^{\lambda(t-t_0)} \geq c_0 > 0$
        (strictly increasing while in Collapse tier).
  \item For $t - t_0 \gg 1/\lambda$: $C(t;t_0) \sim (c_0/\lambda) e^{\lambda(t-t_0)}$
        (exponential regime).
  \item The penalty is bounded by $C(\tau_C;t_0) = \frac{c_0}{\lambda}(e^{\lambda\Delta\tau_C} - 1)$
        where $\Delta\tau_C = \tau_C - t_0$ is the Collapse-tier
        dwell time.
\end{enumerate}
\end{proposition}

\begin{proof}
(i) $C(t_0;t_0) = \frac{c_0}{\lambda}(e^0 - 1) = 0$.
(ii) $\frac{d}{dt}C = c_0 e^{\lambda(t-t_0)} \geq c_0$ since $\lambda, t-t_0 \geq 0$.
(iii)--(iv) follow from the explicit formula. \qed
\end{proof}

\begin{remark}[Runtime Implementation Gap]
No current implementation of $C(t;t_0)$ exists in \pmhq{}.
The closest analog is the telemetry counter
\texttt{pirtm\_drift\_tier\_total\{tier="collapse"\}}, which counts
Collapse-tier events but does not accumulate the integral.
A penalty accumulator module is recommended as a post-ADR-041
deliverable (see Section~8 of the main report).
\end{remark}

% ============================================================
\section{Summary of Numerical Constants}
\label{sec:constants-summary}
% ============================================================

\noindent For reference, all constants used in proofs above with
their ADR-010 source values:

\begin{center}
\renewcommand{\arraystretch}{1.3}
\begin{tabular}{lllp{5.5cm}}
\toprule
\textbf{Symbol} & \textbf{Code Name} & \textbf{Value} & \textbf{Role} \\
\midrule
$B$ & \texttt{OPERATOR\_CONTRACTIVITY\_BOUND} & $1.0$
  & Expansion threshold (Thm.~\ref{thm:l0-contractivity}) \\
$\eps$ & \texttt{DEFAULT\_EPSILON} & $0.05$
  & Contraction margin; $B_\eps = 0.95$ \\
$\eps_C$ & \texttt{DRIFT\_THRESHOLD} & $0.30$
  & Collapse tier threshold (L0-003 policy) \\
$\eta$ & \texttt{MARSHALLING\_FIDELITY\_THRESHOLD} & $0.95$
  & L1 fidelity gate (Thm.~\ref{thm:l1-fidelity}) \\
$\sigma$ & \texttt{SPECTRAL\_CONSERVATION\_TOLERANCE} & $0.02$
  & Eigenvalue tolerance (Thm.~\ref{thm:spectral-conservation}) \\
$\eps_W$ & Watch tier (proposed) & $0.05$
  & $= \eps$; entry threshold for Watch \\
$\eps_X$ & Warn tier (proposed) & $0.10$
  & $= 2\eps$; entry threshold for Warn \\
$\rho_{\wac}$ & \texttt{maxWacAllowed} & $1.05$
  & WAC tolerance (incommensurable, Cor.~\ref{cor:three-way}) \\
\bottomrule
\end{tabular}
\end{center}

\begin{importantbox}
\textbf{Key Numerical Relationships Established in This Appendix:}
\[
  \underbrace{0.02}_{\sigma}
  \;<\;
  \underbrace{0.04}_{2\sigma}
  \;<\;
  \underbrace{0.05}_{\eps = \eps_W}
  \;<\;
  \underbrace{0.10}_{\eps_X}
  \;<\;
  \underbrace{0.30}_{\eps_C = \texttt{DRIFT\_THRESHOLD}}
  \;<\;
  \underbrace{0.95}_{B_\eps = \eta}
  \;<\;
  \underbrace{1.00}_{B}
  \;<\;
  \underbrace{1.05}_{\rho_\wac}
\]
This chain confirms internal consistency: the spectral tolerance
$\sigma$ is smaller than the contraction margin $\eps$
(Corollary~\ref{cor:sigma-eps-consistency}), the Watch/Warn
tier thresholds are well-separated from the Collapse threshold,
and the Collapse threshold lies safely within the stable regime
$[0, B_\eps)$.
\end{importantbox}

% ============================================================
%  END OF APPENDIX
% ============================================================


\nocite{*}
\bibliographystyle{unsrt}  
\bibliography{references}
\end{document}